% ============================================================
%  MTH636 -- Game Theory | Lecture 10
%  Existence Theorems & Infinite Strategy Spaces
% ============================================================
\documentclass[12pt, a4paper]{article}

% ============================================================
%  MTH636 -- Game Theory
%  Notes Preamble — Color-Coded, TikZ-Enhanced
% ============================================================

% --- Core packages ---
\usepackage[T1]{fontenc}
\usepackage[utf8]{inputenc}
\usepackage{lmodern}
\usepackage[margin=2.5cm, top=2.8cm, bottom=2.8cm]{geometry}
\usepackage{microtype}

% --- Math ---
\usepackage{amsmath, amssymb, amsthm, mathtools}
\usepackage{bm}        % bold math
\usepackage{cancel}    % strike-through in math

% --- Color & Boxes ---
\usepackage[dvipsnames, svgnames, x11names, table]{xcolor}
\usepackage[most]{tcolorbox}
\tcbuselibrary{breakable, skins, theorems}

% --- TikZ & Diagrams ---
\usepackage{tikz}
\usepackage{tikz-qtree}
\usetikzlibrary{arrows.meta, positioning, shapes.geometric, fit, backgrounds,
                decorations.pathreplacing, calc, matrix, trees}
\usepackage{pgfplots}
\pgfplotsset{compat=1.18}

% --- Tables ---
\usepackage{booktabs}
\usepackage{multirow}
\usepackage{array}
\usepackage{tabularx}
\usepackage{makecell}

% --- Typography & Layout ---
\usepackage{enumitem}
\usepackage{parskip}
\usepackage{fancyhdr}
\usepackage{titlesec}
\usepackage{titletoc}
\usepackage{hyperref}
\usepackage{graphicx}
\usepackage{float}
\usepackage{caption}

% --- Color Palette ---
\definecolor{defblue}{HTML}{1A5276}      % Definition header
\definecolor{defbluebg}{HTML}{D6EAF8}    % Definition background
\definecolor{thmgreen}{HTML}{1E8449}     % Theorem header
\definecolor{thmbg}{HTML}{D5F5E3}        % Theorem background
\definecolor{exorange}{HTML}{B7770D}     % Example header
\definecolor{exbg}{HTML}{FEF9E7}         % Example background
\definecolor{warnred}{HTML}{922B21}      % Warning/Key Insight header
\definecolor{warnbg}{HTML}{FADBD8}       % Warning background
\definecolor{noticegray}{HTML}{2C3E50}   % Note header
\definecolor{noticebg}{HTML}{EBF5FB}     % Note background
\definecolor{taketawayyellow}{HTML}{7D6608} % Key Takeaway header
\definecolor{takeawaybg}{HTML}{FDFEFE}   % Takeaway background
\definecolor{accentpurple}{HTML}{6C3483} % Accent / proof
\definecolor{lightgray}{HTML}{F2F3F4}    % shading

% ============================================================
%  TCOLORBOX ENVIRONMENTS
% ============================================================

% -- Definition (Blue) --
\newtcolorbox{defbox}[2][]{
  enhanced, breakable,
  colback=defbluebg, colframe=defblue, coltitle=white,
  fonttitle=\bfseries\sffamily,
  title={Definition: #2},
  attach boxed title to top left={yshift=-2mm, xshift=4mm},
  boxed title style={colback=defblue, rounded corners},
  arc=4pt, boxrule=1.2pt,
  left=6pt, right=6pt, top=8pt, bottom=6pt,
  #1
}

% -- Theorem / Proposition (Green) --
\newtcolorbox{thmbox}[2][]{
  enhanced, breakable,
  colback=thmbg, colframe=thmgreen, coltitle=white,
  fonttitle=\bfseries\sffamily,
  title={Theorem: #2},
  attach boxed title to top left={yshift=-2mm, xshift=4mm},
  boxed title style={colback=thmgreen, rounded corners},
  arc=4pt, boxrule=1.2pt,
  left=6pt, right=6pt, top=8pt, bottom=6pt,
  #1
}

% -- Proposition (Green variant) --
\newtcolorbox{propbox}[2][]{
  enhanced, breakable,
  colback=thmbg, colframe=thmgreen!70!black, coltitle=white,
  fonttitle=\bfseries\sffamily,
  title={Proposition: #2},
  attach boxed title to top left={yshift=-2mm, xshift=4mm},
  boxed title style={colback=thmgreen!70!black, rounded corners},
  arc=4pt, boxrule=1.2pt,
  left=6pt, right=6pt, top=8pt, bottom=6pt,
  #1
}

% -- Example (Orange/Yellow) --
\newtcolorbox{exbox}[2][]{
  enhanced, breakable,
  colback=exbg, colframe=exorange, coltitle=white,
  fonttitle=\bfseries\sffamily,
  title={Example: #2},
  attach boxed title to top left={yshift=-2mm, xshift=4mm},
  boxed title style={colback=exorange, rounded corners},
  arc=4pt, boxrule=1.2pt,
  left=6pt, right=6pt, top=8pt, bottom=6pt,
  #1
}

% -- Key Insight / Warning (Red) --
\newtcolorbox{warnbox}[2][]{
  enhanced, breakable,
  colback=warnbg, colframe=warnred, coltitle=white,
  fonttitle=\bfseries\sffamily,
  title={Key Insight: #2},
  attach boxed title to top left={yshift=-2mm, xshift=4mm},
  boxed title style={colback=warnred, rounded corners},
  arc=4pt, boxrule=1.2pt,
  left=6pt, right=6pt, top=8pt, bottom=6pt,
  #1
}

% -- Note (Gray-Blue) --
\newtcolorbox{notebox}[1][]{
  enhanced, breakable,
  colback=noticebg, colframe=noticegray!60, coltitle=white,
  fonttitle=\bfseries\sffamily,
  title={Note},
  attach boxed title to top left={yshift=-2mm, xshift=4mm},
  boxed title style={colback=noticegray!60, rounded corners},
  arc=4pt, boxrule=0.8pt,
  left=6pt, right=6pt, top=8pt, bottom=6pt,
  #1
}

% -- Key Takeaway --
\newtcolorbox{takebox}[1][]{
  enhanced, breakable,
  colback=takeawaybg, colframe=taketawayyellow, coltitle=white,
  fonttitle=\bfseries\sffamily,
  title={Key Takeaway},
  attach boxed title to top left={yshift=-2mm, xshift=4mm},
  boxed title style={colback=taketawayyellow, rounded corners},
  arc=4pt, boxrule=1.2pt,
  left=6pt, right=6pt, top=8pt, bottom=6pt,
  #1
}

% -- Intuition --
\newtcolorbox{intubox}[1][]{
  enhanced, breakable,
  colback=accentpurple!8, colframe=accentpurple!60, coltitle=white,
  fonttitle=\bfseries\sffamily,
  title={Intuition},
  attach boxed title to top left={yshift=-2mm, xshift=4mm},
  boxed title style={colback=accentpurple!80, rounded corners},
  arc=4pt, boxrule=0.8pt,
  left=6pt, right=6pt, top=8pt, bottom=6pt,
  #1
}

% ============================================================
%  PAYOFF MATRIX HELPER
%  Usage: inside tabular, use \payoff{a}{b} for (a,b)
% ============================================================
\newcommand{\payoff}[2]{$(#1,\;#2)$}
\newcommand{\nash}[2]{\underline{$#1$},\underline{$#2$}}  % underline NE payoffs
\newcommand{\Nset}{\mathbb{N}}
\newcommand{\Rset}{\mathbb{R}}
\newcommand{\Sset}{\mathcal{S}}
\newcommand{\Ui}{u_i}

% ============================================================
%  SECTION FORMATTING
% ============================================================
\titleformat{\section}
  {\large\bfseries\sffamily\color{defblue}}
  {\thesection}{1em}{}[\titlerule]

\titleformat{\subsection}
  {\normalsize\bfseries\sffamily\color{thmgreen!80!black}}
  {\thesubsection}{1em}{}

\titleformat{\subsubsection}
  {\normalsize\itshape\sffamily\color{exorange!80!black}}
  {\thesubsubsection}{1em}{}

% ============================================================
%  HEADER / FOOTER
% ============================================================
\pagestyle{fancy}
\fancyhf{}
\fancyhead[L]{\sffamily\small\color{defblue} MTH636 -- Game Theory}
\fancyhead[R]{\sffamily\small\color{defblue} \leftmark}
\fancyfoot[C]{\sffamily\small\thepage}
\renewcommand{\headrulewidth}{0.4pt}

% ============================================================
%  HYPERREF SETUP
% ============================================================
\hypersetup{
  colorlinks=true,
  linkcolor=defblue,
  urlcolor=thmgreen,
  citecolor=accentpurple,
  bookmarksnumbered=true,
  pdftitle={MTH636 Game Theory Notes},
  pdfauthor={MTH636 Student Notes}
}

% ============================================================
%  MISC HELPERS
% ============================================================
\newcommand{\BR}{\text{BR}}        % Best Response
\newcommand{\NE}{\text{NE}}        % Nash Equilibrium
\newcommand{\SPNE}{\text{SPNE}}    % Subgame Perfect NE
\newcommand{\Dom}{\text{dom}}      % Dominance
\newcommand{\mixs}[1]{\sigma_{#1}} % Mixed strategy
\newcommand{\pures}[1]{s_{#1}}     % Pure strategy


% ============================================================
\begin{document}

% Title
\begin{center}
  {\LARGE\bfseries\sffamily\color{defblue} MTH636 -- Game Theory}\\[6pt]
  {\Large\sffamily Lecture 10: Existence of Equilibria}\\[4pt]
\end{center}
\vspace{0.5cm}
\hrule height 1.5pt
\vspace{0.3cm}

% ============================================================
\section{Fundamental Theorems of Game Theory}

Two of the most important theorems in game theory guarantee the existence of solutions, provided that we allow players to use mixed strategies and the game is finite.

\begin{thmbox}{Von Neumann's Minimax Theorem (1928)}
  Every two-player zero-sum game where both players have \textbf{finitely many pure strategies} has a \textbf{value} in mixed strategies.
  
  That is, $\underline{V} = \overline{V}$ when players are allowed to randomize.
\end{thmbox}

\begin{thmbox}{Nash's Existence Theorem (1950)}
  Every strategic form game where all players have \textbf{finitely many pure strategies} has at least one \textbf{Nash Equilibrium} in mixed strategies.
\end{thmbox}

\begin{notebox}
  These theorems require the set of pure strategies to be \textbf{finite}. If the strategy space is infinite (e.g., choosing any real number, or choosing any integer), an equilibrium or value is \textbf{not guaranteed} to exist, even in mixed strategies.
\end{notebox}

% ============================================================
\section{Counterexample: Infinite Strategy Spaces}

Let's see an example where Von Neumann's Minimax Theorem fails because the strategy spaces are infinite.

\begin{exbox}{The Integer Game}
  $N = \{1, 2\}$. Both players independently choose a natural number: $S_1 = S_2 = \mathbb{N} = \{1, 2, 3, \dots\}$.
  The player who picks the larger number wins:
  \[
    u_1(s_1, s_2) = 
    \begin{cases} 
      1 & \text{if } s_1 > s_2 \\
      0 & \text{if } s_1 = s_2 \\
      -1 & \text{if } s_1 < s_2
    \end{cases}
  \]
  This is a zero-sum game.
\end{exbox}

\subsection{Pure Strategy Value}
\begin{itemize}
    \item \textbf{Player 1's Max-Min:} $\sup_{s_1} \inf_{s_2} u_1(s_1, s_2) = -1$. (No matter what $s_1$ Player 1 chooses, P2 can always pick $s_2 = s_1 + 1$ to make P1 lose).
    \item \textbf{Player 2's Min-Max:} $\inf_{s_2} \sup_{s_1} u_1(s_1, s_2) = 1$.
\end{itemize}
Since $-1 \neq 1$, there is no value in pure strategies.

\subsection{Mixed Strategy Value}
Does it have a value in mixed strategies? 
Let $\Sigma_i$ be the set of probability distributions over $\mathbb{N}$.
Fix any mixed strategy $\sigma_1 \in \Sigma_1$ for Player 1. Since probabilities must sum to 1, most of the probability mass must be concentrated on smaller numbers. 
For any $\epsilon > 0$, there exists an integer $k$ such that Player 1 plays a number $\leq k$ with probability $> 1-\epsilon$:
\[
  \sigma_1(\{1, 2, \dots, k\}) > 1 - \epsilon
\]
If Player 2 simply plays the pure strategy $s_2 = k+1$, Player 2 wins against all of P1's numbers from $1$ to $k$, and only potentially loses to numbers $> k$.
\[
  U_1(\sigma_1, s_2 = k+1) = \underbrace{\sigma_1(\{k+2, \dots\})}_{\text{P1 wins}} (1) + \underbrace{\sigma_1(\{1, \dots, k\})}_{\text{P1 loses}} (-1)
\]
Since $\sigma_1(\{1, \dots, k\}) > 1-\epsilon$, the probability P1 wins is less than $\epsilon$.
\[
  U_1(\sigma_1, k+1) < \epsilon(1) + (1-\epsilon)(-1) = -1 + 2\epsilon
\]
Since this is true for any $\epsilon$, the infimum over all of P2's strategies is $\leq -1$.
\[
  \inf_{\sigma_2 \in \Sigma_2} U_1(\sigma_1, \sigma_2) \leq -1 \quad \implies \quad \sup_{\sigma_1} \inf_{\sigma_2} U_1(\sigma_1, \sigma_2) = -1
\]
By symmetry for Player 2, $\inf_{\sigma_2} \sup_{\sigma_1} U_1(\sigma_1, \sigma_2) = 1$.
Thus, \textbf{the value does not exist even in mixed strategies}.

% ============================================================
\section{Simplifying Minimax for Mixed Strategies}

For finite two-player zero-sum games, finding the mixed max-min value involves continuous optimization over the simplex $\Sigma_i$.

\begin{thmbox}{Minimizing over Pure Strategies}
  When calculating the max-min value in mixed strategies, the inner minimization over all mixed strategies $\sigma_2$ can be replaced by a minimization over just the \textbf{pure} strategies $s_2$:
  \[
    \underline{V} = \max_{\sigma_1 \in \Sigma_1} \min_{\sigma_2 \in \Sigma_2} U(\sigma_1, \sigma_2) = \max_{\sigma_1 \in \Sigma_1} \min_{s_2 \in S_2} U(\sigma_1, s_2)
  \]
  Similarly:
  \[
    \overline{V} = \min_{\sigma_2 \in \Sigma_2} \max_{\sigma_1 \in \Sigma_1} U(\sigma_1, \sigma_2) = \min_{\sigma_2 \in \Sigma_2} \max_{s_1 \in S_1} U(s_1, \sigma_2)
  \]
\end{thmbox}

\textbf{Why?} Because $U(\sigma_1, \sigma_2)$ is a weighted average of $U(\sigma_1, s_2)$ for all pure $s_2$. A weighted average can never be strictly less than the absolute minimum of its components. So the worst-case response by an opponent will always just be to put 100\% probability on the single pure strategy that hurts you the most.

\end{document}